% !TEX program = pdflatex
\documentclass[11pt]{article}
\usepackage[T1]{fontenc}
\usepackage{lmodern,amsmath,amssymb,amsthm,booktabs,needspace}
\usepackage[a4paper,margin=27mm]{geometry}
\usepackage[hidelinks]{hyperref}
\DeclareMathOperator{\area}{area}\DeclareMathOperator{\dinv}{dinv}
\DeclareMathOperator{\defc}{defc}\DeclareMathOperator{\sym}{sym}
\newtheorem{theorem}{Theorem}\newtheorem{lemma}[theorem]{Lemma}
\newtheorem{proposition}[theorem]{Proposition}
\newtheorem{claim}{Claim}
\title{On Dyck paths with deficit at most $3n-15$}\author{Graham Hawkes}\date{}
\begin{document}\maketitle
\begin{abstract}
We give a combinatorial formula for the joint distribution of area and dinv on Dyck paths with deficit at most $3n-15$. The formula expresses the corresponding homogeneous components of the $q,t$-Catalan polynomial as nonnegative sums of symmetric, unimodal polynomials, with a combinatorial interpretation of the multiplicities.
\end{abstract}

\section{Introduction}
We use the following terminology and notation from~\cite{H}:
\begin{itemize}
\item From \S1.1: Dyck sequences, affine and reverse affine sequences, extractable elements, Dyck skeletons and full Dyck skeletons, lowers and uppers, area, dinv pairs, and deficit ($\defc$).
\item From \S2: keys, their associated lowers and uppers, and balance on a key.
\item From \S3.2: dual and reverse dual sequences.
\item From \S3.3: Dyck tableaux, their row-reading words, and $pTs$ objects.
\item From \S6.3: the families $\mathcal F_n$, $\mathcal U_n$, and $\mathcal E_n$, the counts $\mathcal N_n(d,a)$, and the polynomials $\sym(i,j)$.
\end{itemize}
We identify a Dyck path of semilength $n$ with its length-$n$ Dyck sequence. Write $|W|$ for word length, $U:V$ for concatenation, and $P[-1]$ for the last entry of $P$. Set $M=\binom n2$ and $C_n(q,t)=\sum_D q^{\area(D)}t^{\dinv(D)}$, summed over the length-$n$ Dyck sequences. 

\begin{theorem}\label{main}
The terms of total degree at least $M-3n+15$ in $C_n(q,t)$ are given by
\[
 \left.C_n(q,t)\right|_{\deg_{q,t}\ge M-3n+15}
 =\sum_{d=0}^{3n-15}\sum_{a=0}^{d}\mathcal N_n(d,a)\,\sym(a,M-d-a).
\]
\end{theorem}
Let $c_d(a)$ count the paths of deficit $d$ and area $a$, with $c_d(-1)=0$.
\begin{lemma}\label{first-difference}
For $0\le a\le d\le3n-15$,
\[
 \mathcal N_n(d,a)=c_d(a)-c_d(a-1).
\]
\end{lemma}
Proving Lemma~\ref{first-difference} suffices to prove Theorem~\ref{main}. Summing the identity determines the coefficients for areas $0\le a\le d$. The flat-middle theorem~\cite[Theorem~3.12]{H} gives the constant coefficient $c_d(d)$ for $d\le a\le M-2d$, and the symmetry $C_n(q,t)=C_n(t,q)$~\cite[\S1]{LLL} determines the remaining coefficients by reflection $a\mapsto M-d-a$.

\section{Adjacent area counts}
For a key $K:=(P,X,h)$, put $m=\max(P)$ and write $|X|=2k+1$, so $n=|P|+|X|$. Write $J(P,X)$ for the number of pairs $(v,x)\in P\times X$ with $v-x\notin\{0,1\}$. In this notation, the formulas in~\cite[\S2]{H} give
\begin{equation}\label{stats}
 \begin{aligned}
 \defc(K)&=\defc(P)+J(P,X)+\binom{|X|}{2}-h-\sum_{x\in X}x,\\
 \area(K)&=\area(P)+\sum_{x\in X}x+k.
 \end{aligned}
\end{equation}
Here $\defc(K)$ is the deficit of any lower or upper on $K$, and $\area(K)$ is the area of any lower on $K$. The area of an upper is $\area(K)+1$. For fixed $n$, let $\kappa(d,a)$ denote the set of keys $K=(P,X,h)$ with $|P|+|X|=n$, $\defc(K)=d$, and $\area(K)=a$. The statistic-preserving bijections~\cite[\S1.2]{H} identify all paths with lowers and all paths other than full skeletons with uppers. Hence
\begin{equation}\label{diff}
 \begin{aligned}
 c_d(a)-c_d(a-1)&=\#\{F\in\mathcal F_n:\defc(F)=d,\ \area(F)=a\}\\
 &\quad+\sum_{K\in\kappa(d,a-1)}\bigl(\#\mathrm{uppers}(K)-\#\mathrm{lowers}(K)\bigr).
 \end{aligned}
\end{equation}
Here $\#\mathrm{uppers}(K)$ and $\#\mathrm{lowers}(K)$ count the uppers and lowers associated to $K$, respectively.

The next lemma bounds $h$ and the multiplicities in $X$ for an \emph{unbalanced} key (one with unequal numbers of uppers and lowers). Let $x_v$ denote the multiplicity of $v$ in $X$.
\begin{lemma}\label{bounds}
For an unbalanced key,
\[
x_v+x_{v-1}\le k+1,\qquad
 x_v+x_{v-1}\le k\quad\text{if }v\le\min(P[-1]+1,m), \qquad   h\le k^2,.
\]
Moreover $\max X\ge\min(P[-1]+2,m+1)$.
\end{lemma}
\begin{proof}
We use the $pTs$ objects of~\cite[\S3.3]{H}, with the definitions adjusted to all $k\ge0$ and with endpoints $P[-1]$ and $m-1$ in place of $0$ and $-1$. A nonempty corner term in the telescope supplies a $pTs$ object with outer lengths $(k+1,0)$ or $(0,k+1)$ and tableau column lengths $(k,0)$. If both corner terms vanish, the telescoping and tableau arguments of~\cite[\S3.2, Lemma~3.4, and Corollary~3.5]{H} express the nonzero imbalance on $K$ as a signed sum over $pTs$ objects. In either case, there exists a $pTs$ object $(p,T,s)$ with content $X$ and $\dinv(p:\mathrm{RR}(T):s)=h$. We have $|p|\ne|s|$, and the column lengths of $T$ are $k-\min(|p|,|s|)$ and $k+1-\max(|p|,|s|)$.

The factors $p$, the rows of $T$, and $s^{rev}$ form dual Dyck sequences.  By the restrictions above at most $k+1$ are non-empty and since each contains at most one of $x_v$ or $x_{v+1}$ the first assertion follows. If the hypothesis in the second assertion holds neither $p$ nor $s$ contain either of these elements so the second concusion holds. Since the factors we are considering are dual Dyck they have no internal $\dinv$ and the $\dinv$ of the $pTs$ object is just the sum over all pairs of factors of the $\dinv$ between the two factors, which is at most the minumum length of the the two factors.  So we have a factorization of $2k+1$ elements into $k+1$ parts (some may be empty) and will to consider, over such factorizations, what the maximum of the possible sum over all pairs of the minimum factor length in that pair is.  But, given $k+1$ cups and $2k+1$ coins it is an elementary combinatorics problem to show the maximum sum over pairs of cups of the number of coins in the less full cup is $k^2$.  

Finally, a nonempty $p$ contains an entry at least $P[-1]+2$, and a nonempty $s$ contains an entry at least $m+1$. Since at least one is nonempty, the bound on $\max X$ follows.

\end{proof}






\section{Classification of the unbalanced keys}
Throughout this section, let $K=(P,X,h)$ be an unbalanced key with $\area(K)<\defc(K)\le3n-15$, and let $(P,A,B)$ be an associated lower or upper. Define
\[
 \begin{aligned}
 \xi&=\defc(K)-(3n-15)\le0,\\
 \zeta&=\defc(K)+\area(K)-2(3n-15)<0.
 \end{aligned}
\]
\begin{claim}\label{k-bound}
$k\le5$.
\end{claim}
\begin{proof}
By Lemma~\ref{bounds}, each entry of $P$ forms at least $k$ pairs counted by $J(P,X)$, so $J(P,X)\ge k|P|$. Using this bound, $h\le k^2$ from the same lemma, and~\eqref{stats}, and dropping the nonnegative term $\defc(P)+\area(P)$ gives
\[
 \zeta\ge k|P|+\binom{2k+1}{2}-k^2+k-6n+30=(k-6)(|P|+k-4),
\]
where we have used $n=|P|+2k+1$ to derive the equality. If $k\ge6$, the right-hand side is nonnegative, contradicting $\zeta<0$. Thus $k\le5$.
\end{proof}

Let $G$ be the submultiset of $X$ consisting of entries greater than $m$, and write
\[
 R=\sum_{x\in G}(x-(m+1)).
\]
We first bound $R$ for use in the classification.

\begin{claim}\label{excess-bounds}
We have $R\le\binom{|G|}{2}$. If both $A$ and $B$ contain entries above $m$, then $R\le\binom{|G|-1}2$. If $A$ consists of $k+1$ entries above $m$ and all entries of $B$ are at most $m$, then
\begin{equation}\label{oneword}
 h\le(k^2-R)+x_mx_{m+1}.
\end{equation}
\end{claim}
\begin{proof}
If $G$ is empty, then $R=0$. Otherwise, the endpoint conditions on $A$ and $B$ imply that every integer from $m+1$ to $\max(G)$ occurs in $G$. Thus $R \le0+1+\cdots+(|G|-1)=\binom{|G|}{2}$.

If both $A$ and $B$ contain an entry greater than $m$, then each contains $m+1$. Thus $m+1$ occurs at least twice in $G$, and so
 $R\le0+0+1+\cdots+(|G|-2)=\binom{|G|-1}2$.

For~\eqref{oneword}, there are $\binom{k+1}2+\binom k2=k^2$ pairs internal to $A$ or $B$. Each occurrence of $x>m+1$ in $A$ has earlier entries $m+1,\ldots,x-1$. These give $R$ distinct internal pairs that cannot be dinv pairs, so $\dinv(A)+\dinv(B)\le k^2-R$. Let $C(A,B)$ count the dinv pairs between $A$ and $B$. Every such pair has the form $(m+1,m)$, so $C(A,B)\le x_mx_{m+1}$. Adding the internal and cross contributions gives
\[
 h=\dinv(A)+\dinv(B)+C(A,B)\le(k^2-R)+x_mx_{m+1}.
\]
\end{proof}


\Needspace{6\baselineskip}
\begin{claim}\label{below-maximum}
If $P[-1]<m$, then $k=0$, $m=2$, and $P[-1]=0$. The key has exactly one associated lower or upper, which is an upper in $\mathcal U_n$.
\end{claim}
\begin{proof}
Since $P[-1]$ is not extractable, $P$ is a full Dyck skeleton. By~\cite[Lemma~1.12]{H}, every positive entry has at least two entries of each value smaller than it to its left. Consider the the subword of the first leftmost two entries of each element less than $m$, the first $m$, and $P[-1]$:
\[
 Q=00\,11\cdots(m-1)(m-1)\,m\,P[-1].
\]
The deficit of $P$ is the number of non-dinv pairs minus the area. This can be broken two summands: the non-dinv pairs within $Q$ minus the area of $Q$ plus the non-dinv pairs with at least one element in $P \setminus Q$ minus the area of $P \setminus Q$.  Without its last entry, it is easy to count that the non-dinv pairs minus the area of  $Q \setminus Q[-1]$ is $m^2$. Appending $P[-1]$ adds at least $m-1$ (handle $P[-1]=m-1$ and $P[-1]<m-1$ separately). Thus the first in the two summands above is at least $m^2+m-1$.
For each $v$ in $P \setminus Q$ the contribution to the second summand is at least $\max(v,m-1)+\mathbf{1}_{\{v-P[-1]\notin\{0,1\}\}}$ (handle $v=m$, $v=m-1$, $v<m-1$ separately). Thus
\begin{equation}\label{marked-deficit}
 \defc(P)\ge m^2+m-1+\sum_{v\in P\setminus Q}\left(\max(v,m-1)+\mathbf{1}_{\{v-P[-1]\notin\{0,1\}\}}\right).
\end{equation}

Suppose $k\ge1$. In this case we drop the indicator term in~\eqref{marked-deficit}. By Lemma~\ref{bounds}, each entry $v \in P \setminus Q$ forms $2k+1-x_v-x_{v-1}\ge k$ non-dinv pairs counted by $J(P\setminus Q,X)$, with the bound improving to $k+1$ when $v\le\min(P[-1]+1,m)$. Now consider $J(P\setminus Q,X)+\max(v,m-1)$: It gives at least $1+2$ if $m\ge3$. If $m=2$, the improved bound applies to $v=0,1$, giving at least $(k+1)+1\ge3$, while $v=2$ gives at least $k+2\ge3$. If $m=1$, then $P[-1]=0$ and $x_0+x_1\le k$, so $|G|\ge k+1$. Each factor reaching above $1$ must contain a $1$ by the endpoint conditions. Since the factor lengths are $k+1$ and $k$, neither factor can contain all of $G$; hence both contain a $1$, giving $x_1\ge2$ and therefore $k\ge2$. The improved bound applies to both possible values $v=0,1$, so $J(P\setminus Q,X)  \ge k+1\ge3$. Therefore:
\[
 \sum_{v\in P\setminus Q}\max(v,m-1)+J(P\setminus Q,X)\ge3|P\setminus Q|,
\]
and consequently $\defc(P)+J(P\setminus Q,X)\ge m^2+m-1+3|P\setminus Q|$.

We next bound $J(Q,X)-\sum X$ by counting the contribution of each entry $x\in X$. Since $|Q|=2m+2$, this contribution is $2m+2$ minus the number of entries of $Q$ equal to $x$ or $x+1$, minus $x$. If $x \le m-1$ (handle $x=m-1$ and $x<m-1$ separately) this is at least $m-1$. For $x=m$ this is $m+1$. For $x>m$, it is $2m+2-x=m+1-(x-m-1)$. Summing over $X$ gives
\[
 J(Q,X)-\sum X\ge(m-1)|X|+2(x_m+|G|)-R.
\]

Using the bound $h\le k^2$ in~\eqref{stats}, we get
\[
 \begin{aligned}
 \defc(K)\ge{}&m^2+m-1+3|P\setminus Q|\\
 &+(m-1)|X|+2(x_m+|G|)-R+k(k+1).
 \end{aligned}
\]

Each factor reaching above $m$ contains an entry $m$, so the factors of lengths $k+1$ and $k$ contain at most $k$ and $k-1$ entries above $m$, so
\[
 R\le\binom k2+\binom{k-1}2=(k-1)^2.
\]
Using $|X|=2k+1$ and $n=|P\setminus Q|+2m+2+|X|$, the preceding bound on $\defc(K)$ therefore gives
\[
 \defc(K)-(3n-15)\ge(m-1)(m-3)+(2m-5)k+2(x_m+|G|).
\]
If $m\ge3$, the right-hand side is at least $k>0$. If $m=2$, Lemma~\ref{bounds} gives $x_0+x_1\le k$, so $x_m+|G|\ge k+1$ and the right-hand side is at least $-1-k+2(k+1)=k+1>0$. If $m=1$, we already have $x_1\ge2$ and $|G|\ge k+1$, so the right-hand side is at least $-3k+2(k+3)=6-k>0$ by Claim~\ref{k-bound}. In every case, $\defc(K)>3n-15$, contradicting the hypothesis.

Thus we have $k=0$, $X:=\{x\}$, $h=0$, and~\eqref{stats} gives $\defc(K)=\defc(P)+J(P,\{x\})-x$. The only way for the key to be imbalanced is if no lower and one upper exists and $P[-1]+2\le x\le m$, so $m\ge2$.
Suppose $m\ge3$. Then adding $J(P\setminus Q,\{x\})$ to both sides of ~\eqref{marked-deficit} gives
\[
 \defc(P)+J(P\setminus Q,\{x\})\ge m^2+m-1+3|P\setminus Q|,
\]
since no $v \in P \setminus Q$ can form a $\dinv$ pair with both $P[-1]$ and $x$. 
Since $J(Q,\{x\})-x\ge m-1$ and $n=|P\setminus Q|+2m+3$, we obtain
\[
 \begin{aligned}
 \defc(K)&\ge m^2+m-1+3|P\setminus Q|+(m-1)\\
 &=3n-15+(m-2)^2>3n-15,
 \end{aligned}
\]
contradicting the hypothesis. Therefore $m=2$, and $P[-1]+2\le x\le m$ forces $P[-1]=0$ and $x=2$, giving the unique upper in $\mathcal U_n$.
\end{proof}

\Needspace{20\baselineskip}
For the remaining claims, suppose $P[-1]=m$ and put $r=|P|-m-1$.

\begin{claim}\label{maximum-estimates}
We have $|G|\ge1$. If $k=0$, the key has exactly one associated lower or upper: a lower with $(A,B)=((m+1),\varnothing)$ and
\[
 \defc(K)\ge(m+1)r.
\]
For $k\ge1$,
\begin{align}
 \defc(K)&\ge \binom{|X|}{2}-h-R+\bigl(m+\max(k+1,|G|)\bigr)r,\label{D}\\
 \defc(K)+\area(K)&\ge \binom{m+1}2+(m-1)|X|+x_m+2|G|+k(k+2)\notag\\
             &\qquad+\bigl(m+\max(k+1,|G|)\bigr)r.\label{G}
\end{align}
In particular, if $k=1$, then $\defc(K)\ge1+(m+2)r$.
\end{claim}
\begin{proof}
Mark the \emph{last} occurrence of each value in $P$. In order these form $S_m=012\cdots m$, since $P$ ends at its maximum. There are $r$ unmarked entries. Each unmarked entry $v$ has two earlier entries of each smaller value: otherwise the consecutive rise described above would prevent $v$ from repeating. It also precedes all $m-v$ larger marked entries. These give $2v+(m-v)$ pairs that are not dinv pairs. In the first set of pairs the larger entry is unmarked; in the second it is marked, so no pair is counted twice. Subtracting the contribution $v$ to area leaves at least $m$ per unmarked entry. Since $\defc(S_m)=0$, we obtain $\defc(P)\ge mr$.

Lemma~\ref{bounds} gives $|G|\ge1$. For $k=0$ the only unbalanced letter is $m+1$, giving one lower and no upper. Equation~\eqref{stats} then gives $\defc(K)\ge(m+1)r$.

For $k\ge1$ every unmarked entry forms at least $\max(k+1,|G|)$ pairs counted by $J(P,X)$. Against $S_m$, subtract $x$ from the number of pairs counted by $J(S_m,(x))$. The result is nonnegative for $x\le m$, since the $x$ smaller levels already contribute $x$ pairs; for $x>m$ it is $-(x-m-1)$. Also $J(S_m,X)=(m-1)|X|+x_m+2|G|$: a letter below $m$ forms two dinv pairs with $S_m$, a letter $m$ one, and a larger letter none. Substitution in~\eqref{stats} yields~\eqref{D} and~\eqref{G}.

For $k=1$, $R\le1$ because only the longer factor can reach $m+2$. Combining this with $h\le1$ in~\eqref{D} gives the last bound.
\end{proof}

\begin{claim}\label{maximum-height}
If $P[-1]=m$, then $m\le1$.
\end{claim}
\begin{proof}
Suppose $m\ge2$. By Claim~\ref{k-bound}, it suffices to consider $0\le k\le5$.

For $k=0$, Claim~\ref{maximum-estimates} and $\defc(K)\le3n-15$ give $(m-2)r\le3m-9$. This excludes $m=2$; for $m\ge3$ it gives $r\le2$, whereas $\area(K)\ge\binom{m+2}2>3(m+4)-15\ge3n-15$, a contradiction.

For $k=1$, the last bound in Claim~\ref{maximum-estimates} gives $(m-1)r\le3m-4$, hence $r\le2$. Take the $pTs$ object supplied by the proof of Lemma~\ref{bounds}. The associated lower or upper has factor lengths $1$ and $2$, so its endpoint conditions give $\max X\le m+2$. An outer word $p$ or $s$ of length at least two would contain an entry at least $m+3$. Since the $pTs$ object also has content $X$, this is impossible. Thus its tableau consists of one dual row of length two, and its remaining entry lies in an outer word and is at least $m+1$. By~\eqref{stats}, $\area(K)\ge\binom{m+2}2+3\ge3m+3\ge3n-15$, again a contradiction.

Now let $2\le k\le5$. If $m\ge3$, the coefficient of $r$ in the lower bound for $\zeta$ obtained from~\eqref{G} is at least $m+k-5\ge0$. After using $x_m+2|G|\ge2$, the constant term in this bound is $(k-3)^2+1$ at $m=3$ and increases by $m+1+|X|-6$ when $m$ increases by one. Thus the bound is positive, contradicting $\zeta<0$.

For $m=2$, Lemma~\ref{bounds} gives $x_0+x_1\le k$, so $x_m+|G|\ge k+1$. Equation~\eqref{G} gives $\zeta\ge(k-3)(k-4)+(k-3)r$. For $k=3,4,5$ this contradicts $\zeta<0$. At $k=2$, the same equation gives
\[
 \zeta\ge-2+x_m+2|G|+(\max(3,|G|)-4)r.
\]
For $|G|\ge4$ the right-hand side is positive, contradicting $\zeta<0$. For $1\le |G|\le3$,~\eqref{D} and Claim~\ref{excess-bounds} give $2r\le3+\binom{|G|}{2}$, hence $r\le |G|$. Since $x_m+|G|\ge3$, we again obtain $\zeta\ge x_m+|G|-2>0$, a contradiction.
\end{proof}

\begin{claim}\label{binary-maximum}
If $P[-1]=m=1$, then $k=0$, and the key has exactly one associated lower or upper: a lower with $(A,B)=((2),\varnothing)$.
\end{claim}
\begin{proof}
For $k=0$, Claim~\ref{maximum-estimates} gives the stated lower. For $k=1$, we have $n=r+5$, so the last bound in that claim gives $\xi\ge1$, contradicting $\xi\le0$. It remains to exclude $2\le k\le5$.

At most $k$ entries are $0$ or $1$, so $|G|\ge k+1$. An unmarked zero forms $|G|+x_1$ pairs counted by $J(P,X)$; an unmarked one forms $|G|$ and contributes one to area. Keeping this distinction in the derivation of~\eqref{G}, and using~\eqref{D}, gives
\begin{align*}
 \zeta&\ge (k-3)(k-5)+x_1+2(|G|-k-1)+(|G|-5+\min(x_1,1))r,\\
 \xi&\ge (k-2)(k-3)+k^2-h-R+(|G|-2)r.
\end{align*}
If $|G|\ge k+2$, the factor $B$ must reach above $1$, so $x_1\ge1$. Both the constant term and the coefficient of $r$ in the first bound are then nonnegative, contradicting $\zeta<0$.

It remains to consider $|G|=k+1$. If $x_1\ge1$, the first bound excludes $k=3,4,5$, since it gives $\zeta\ge0$. At $k=2$ the second bound and Claim~\ref{excess-bounds} give $r\le R\le3$, and the first gives $\zeta\ge3+x_1-r>0$, again contradicting $\zeta<0$.

Suppose now that $|G|=k+1$ and $x_1=0$. All $k+1$ entries above $1$ must lie in $A$ and all $k$ zeros in $B$. Hence $h+R\le k^2$ by~\eqref{oneword}. The bound on $\xi$, together with $\xi\le0$, forces $k=2,3$ and $r=0$. Then $\zeta\ge(k-3)(k-5)\ge0$, contradicting $\zeta<0$.
\end{proof}

\begin{claim}\label{zero-maximum}
If $m=0$, then $k=0$ or $1$, and the key has exactly one associated lower or upper: a lower with
\[
 (A,B)=((1),\varnothing)\quad\text{or}\quad((12),(0)),
\]
respectively.
\end{claim}
\begin{proof}
Here $P$ consists of zeros. For $k=0$, Claim~\ref{maximum-estimates} gives the stated lower. For $k=1$, write $A:B=xyz$. The endpoint conditions give $x\le1$ and $z=0$. The lower condition $y\le x+1$ differs from the upper condition $y\le1$ only for $xyz=120$, giving a single lower: $120$ is the only ordering of $0,1,2$ with the required first and last entries.

It remains to exclude $2\le k\le5$. By Lemma~\ref{bounds}, $x_0\le k$, hence $|G|\ge k+1$. Counting the $|G|\,|P|$ pairs between $P$ and positive entries gives
\begin{align*}
 \xi&=(|G|-3)|P|+\binom{|X|}{2}-h-|G|-R-3|X|+15,\\
 \zeta&\ge(|G|-6)|P|+(k-4)(k-6).
\end{align*}
If $|G|=k+1$, then $x_0=k$, so the bound $x_0+x_1\le k+1$ gives $x_1\le1$. Each factor containing a positive entry must contain a $1$; hence only one factor contains positive entries. All $k+1$ of them must lie in $A$, since $B$ must end in zero. Thus~\eqref{oneword} gives $h+R\le k^2+k$, and so $\xi\ge(k-2)|P|+k^2-7k+11$. For $2\le k\le5$ this contradicts $\xi\le0$ except possibly at $k=3$, $|P|=1$, when $\zeta\ge1$ contradicts $\zeta<0$.

If $|G|\ge k+2$, both factors contain positive entries. For $k=4,5$ we have $\zeta\ge(k-4)(|P|+k-6)\ge0$, contradicting $\zeta<0$. For $k=2,3$, the bound on $\zeta$ excludes $|G|\ge6$. Also $B$ contains a zero, so $|G|\le2k$. These bounds force $|G|=k+2$, and $\zeta<0$ then requires $|P|\ge7-k$. Now Claim~\ref{excess-bounds} gives $R\le\binom{|G|-1}2=\binom{k+1}2$. The displayed bound for $\xi$ gives $\xi\ge |P|-1$ at $k=2$ and $\xi\ge2|P|-5$ at $k=3$, both positive and contrary to $\xi\le0$.
\end{proof}

\Needspace{11\baselineskip}
\begin{proposition}\label{exceptions}
Suppose $K=(P,X,h)$ is an unbalanced key with $\area(K)<\defc(K)\le3n-15$. Exactly one lower or upper $(P,A,B)$ is associated to $K$. If it is an upper, it belongs to $\mathcal U_n$. For a lower, either $P$ consists of zeros, with $(A,B)\in\{((1),\varnothing),((12),(0))\}$; or $P$ is binary and ends in $1$, with $(A,B)=((2),\varnothing)$.
\end{proposition}
\begin{proof}
If $P[-1]<m$, Claim~\ref{below-maximum} gives a single upper in $\mathcal U_n$. If $P[-1]=m$, Claim~\ref{maximum-height} gives $m\le1$. Claims~\ref{binary-maximum} and~\ref{zero-maximum} then give exactly the stated lowers, each unique on its key.
\end{proof}

\section{The exceptional words}
It remains to describe the binary skeleton $P$ in the last case of Proposition~\ref{exceptions}, where $(A,B)=((2),\varnothing)$. It begins with $00$: an initial $01$ would force $P=01$, which violates $\defc(K)\le3n-15$. If $P$ ends in $01$, it has the form $00B01$. Otherwise it ends in $11$; we show that it has the form $0011B11$. Here, as in~\cite[\S6.3]{H}, $B$ in a displayed word form denotes an arbitrary binary word.

Suppose $P$ ends in $11$, and write $z,o$ for its numbers of zeros and ones. Each one contributes to the deficit the number of preceding zeros other than the first. In particular the final two ones contribute $2(z-1)$. Remove the initial $00$ and final $11$. Unless the remaining word starts with $11$, all but at most one of its ones have a zero to their left in that word, hence at least three zeros to their left in $P$. It follows that
\[
 \defc(P)\ge 2(z-1)+2(o-2)-1=2|P|-7.
\]
On the other hand,~\eqref{stats} gives $\defc(K)=\defc(P)+|P|-2\le3(|P|+1)-15$, so $\defc(P)\le2|P|-10$. This contradiction shows that the remaining word starts with $11$, as required.

Pair the lowers listed in Proposition~\ref{exceptions} with full skeletons using the map displayed in~\cite[\S6.3]{H}:
\[
\begin{array}{@{}ll@{}}\toprule
 \text{lower}&\text{full skeleton}\\\midrule
 (0^{n-1},(1),\varnothing)&0010^{n-4}1\\
 (0^{n-3},(12),(0))&001120^{n-6}1\\
 (00B01,(2),\varnothing)&001B12\\
 (0011B11,(2),\varnothing)&00112B12\\\bottomrule
\end{array}
\]
Each image has two zeros before its first one and, when a two occurs, two ones before its first two. The four image families are disjoint: the first two end in one and differ in whether a two occurs; the others end in two and differ in the number of twos. In the last two rows, removing the specified prefix and suffix recovers $B$. Restricting each row to words of length $n$, the table therefore gives a bijection onto $\mathcal E_n$.

In each row the area increases by one and the deficit is unchanged. In the first row the areas are $1,2$ and both deficits are $n-2$; in the second the areas are $4,5$ and both deficits are $2n-7$, by counting pairs of entries. The second row can occur in our range only for $n\ge8$.

For the last two rows, call the entries outside the displayed binary word $B$ the fixed entries. The dinv pairs internal to $B$ do not change. In the third row, each zero in $B$ forms three dinv pairs with the fixed entries, and each one forms two, both before and after the pairing. In the fourth row, each forms four. Among the fixed entries, the number of dinv pairs decreases from three to two in the third row and from seven to six in the fourth. Thus dinv decreases by one as area increases by one, preserving deficit.

Each lower in the table is the unique lower or upper on its key. In the rows with $k=0$, the single letter exceeds $m$, so it cannot be the entry of $B$ in an upper. In the second row, $120$ is the unique admissible ordering. The uppers in $\mathcal U_n$ likewise have no lower, since a $2$ cannot follow $P[-1]=0$ in an affine word.

At fixed deficit $d$, the pairing cancels each negative term from a lower of area $a-1$ in~\eqref{diff} with one full skeleton in $\mathcal E_n$ of area $a$. The remaining terms count the uppers in $\mathcal U_n$ and the full skeletons in $\mathcal F_n\setminus\mathcal E_n$, namely $\mathcal N_n(d,a)$. We obtain
\[
 c_d(a)-c_d(a-1)=\mathcal N_n(d,a)\qquad(0\le a\le d\le3n-15),
\]
including $a=0$ since no lower area is negative. This proves Lemma~\ref{first-difference}.

\begin{thebibliography}{9}
\bibitem{DS} G. Hawkes, \emph{Dyck symmetric functions and applications to $q,t$-Catalan polynomials}, \href{https://arxiv.org/abs/2605.13003v1}{arXiv:2605.13003v1} (2026).
\bibitem{H} G. Hawkes, \emph{An involution on Dyck paths in the region $\defc\le\min(\area,\dinv)$ that interchanges area and dinv}, \href{https://arxiv.org/abs/2609.21403v3}{arXiv:2609.21403v3} (2026).
\bibitem{LLL} K. Lee, L. Li and N.~A. Loehr, \emph{A combinatorial approach to the symmetry of $q,t$-Catalan numbers}, SIAM J. Discrete Math. \textbf{32} (2018), 191--232; \href{https://arxiv.org/abs/1602.01126}{arXiv:1602.01126}.
\end{thebibliography}
\end{document}
